% IEEE Conference Paper Template
% Based on IEEEtran class
% Compile with: pdflatex main.tex
%               bibtex main
%               pdflatex main.tex
%               pdflatex main.tex

\documentclass[conference]{IEEEtran}

% -------------------------------------------------------
% Packages
% -------------------------------------------------------
\usepackage{cite}
\usepackage{amsmath,amssymb,amsfonts}
\usepackage{algorithmic}
\usepackage{graphicx}
\usepackage{textcomp}
\usepackage{xcolor}
\usepackage{hyperref}
\usepackage{booktabs}

\def\BibTeX{{\rm B\kern-.05em{\sc i\kern-.025em b}\kern-.08em
    T\kern-.1667em\lower.7ex\hbox{E}\kern-.125emX}}

% -------------------------------------------------------
% Document
% -------------------------------------------------------
\begin{document}

\title{Self-Organizing Maps for Cart-Pole Control:\\
A Biologically-Inspired Approach to Reinforcement Learning}

\author{
  \IEEEauthorblockN{Author Name}
  \IEEEauthorblockA{
    \textit{Department / Institution} \\
    City, Country \\
    email@institution.edu
  }
  % Uncomment to add a second author:
  % \and
  % \IEEEauthorblockN{Second Author}
  % \IEEEauthorblockA{
  %   \textit{Department / Institution} \\
  %   City, Country \\
  %   email@institution.edu
  % }
}

\maketitle

% -------------------------------------------------------
% Abstract
% -------------------------------------------------------
\begin{abstract}
We present a biologically-inspired control architecture that uses a
Self-Organizing Map (SOM) to solve the classic cart-pole balancing problem.
The SOM discretizes the continuous four-dimensional state space into a
topology-preserving one-dimensional map, enabling efficient policy learning
without deep neural networks.
Experimental results demonstrate that the agent learns a stable balancing
policy within \(N\) episodes, outperforming baseline tabular methods while
remaining interpretable.
\end{abstract}

\begin{IEEEkeywords}
self-organizing map, cart-pole, reinforcement learning,
biologically-inspired control, state discretization
\end{IEEEkeywords}

% -------------------------------------------------------
% I. Introduction
% -------------------------------------------------------
\section{Introduction}
\label{sec:intro}

The cart-pole balancing problem~\cite{barto1983neuronlike} is a canonical
benchmark in control theory and reinforcement learning (RL).
A pole is hinged to a cart that moves along a frictionless track; the agent
must apply left or right forces to keep the pole upright indefinitely.

Classical solutions rely on linear approximations or hand-engineered features,
while modern deep RL approaches~\cite{mnih2015human} achieve strong performance
but lack interpretability.
Self-Organizing Maps~\cite{kohonen1990self} offer an intermediate solution:
they adapt continuously to the input distribution and produce a
topology-preserving discrete representation that supports lightweight tabular
policies.

The remainder of this paper is organized as follows.
Section~\ref{sec:background} reviews relevant background.
Section~\ref{sec:method} describes the proposed architecture.
Section~\ref{sec:experiments} reports experimental results.
Section~\ref{sec:conclusion} concludes and outlines future work.

% -------------------------------------------------------
% II. Background
% -------------------------------------------------------
\section{Background}
\label{sec:background}

\subsection{Self-Organizing Maps}

A Self-Organizing Map~\cite{kohonen1990self} is an unsupervised neural network
that projects high-dimensional input data onto a lower-dimensional (typically
1-D or 2-D) lattice of prototype neurons.
During training, the Best Matching Unit (BMU) and its topological neighbors
are updated toward the current input via:
%
\begin{equation}
  \mathbf{w}_i(t+1) = \mathbf{w}_i(t)
    + \eta(t)\, h(i, \text{BMU}, t)\,
      \bigl[\mathbf{x}(t) - \mathbf{w}_i(t)\bigr],
  \label{eq:som_update}
\end{equation}
%
where \(\eta(t)\) is the learning rate and \(h(i, \text{BMU}, t)\) is a
Gaussian neighborhood kernel that shrinks over time.

\subsection{Cart-Pole Environment}

The cart-pole environment is defined by a four-dimensional continuous state
\(\mathbf{s} = (x,\, \dot{x},\, \theta,\, \dot{\theta})\) and a binary action
space \(\mathcal{A} = \{0, 1\}\) (push left, push right).
An episode ends when \(|\theta| > 12°\), \(|x| > 2.4\,\text{m}\), or 500
time steps are reached.

% -------------------------------------------------------
% III. Method
% -------------------------------------------------------
\section{Proposed Method}
\label{sec:method}

\subsection{Architecture Overview}

Figure~\ref{fig:architecture} illustrates the overall architecture.
Raw state observations are fed into a 1-D SOM encoder that maps them to a
discrete neuron index.
A tabular Q-value table indexed by \((\text{SOM index},\, \text{action})\)
stores learned action values, which are updated by a standard
Q-learning rule~\cite{watkins1992q}.

\begin{figure}[t]
  \centering
  \includegraphics[width=\columnwidth]{figures/architecture.pdf}
  \caption{Overall system architecture. The SOM encodes continuous states
           into discrete indices used by the tabular Q-learner.}
  \label{fig:architecture}
\end{figure}

\subsection{SOM Encoder}

We use a 1-D SOM with \(K\) neurons whose weights \(\mathbf{w}_i \in \mathbb{R}^4\)
are initialized uniformly at random.
The state is normalized to \([0,1]^4\) before computing BMU distances.
Training follows Eq.~\eqref{eq:som_update} with an exponential decay schedule:
%
\begin{align}
  \eta(t) &= \eta_0 \exp\!\left(-\frac{t}{\tau_\eta}\right), \\
  \sigma(t) &= \sigma_0 \exp\!\left(-\frac{t}{\tau_\sigma}\right),
\end{align}
%
where \(\sigma(t)\) governs the neighborhood kernel width.

\subsection{Q-Learning Policy}

Given BMU index \(k\), the agent selects actions via an \(\epsilon\)-greedy
policy and updates Q-values as:
%
\begin{equation}
  Q(k, a) \leftarrow Q(k, a)
    + \alpha\!\left[r + \gamma \max_{a'} Q(k', a') - Q(k, a)\right],
\end{equation}
%
where \(\alpha\) is the learning rate, \(\gamma\) the discount factor,
and \(k'\) the BMU of the next state.

% -------------------------------------------------------
% IV. Experiments
% -------------------------------------------------------
\section{Experiments}
\label{sec:experiments}

\subsection{Setup}

All experiments were conducted on the \texttt{CartPole-v1} environment from
OpenAI Gymnasium~\cite{towers2024gymnasium}.
Hyperparameters are listed in Table~\ref{tab:hyperparams}.

\begin{table}[t]
  \centering
  \caption{Hyperparameter Settings}
  \label{tab:hyperparams}
  \begin{tabular}{lcc}
    \toprule
    Parameter & Symbol & Value \\
    \midrule
    SOM neurons           & \(K\)             & 100   \\
    Initial learning rate & \(\eta_0\)        & 0.5   \\
    Initial neighborhood  & \(\sigma_0\)      & 10    \\
    Q-learning rate       & \(\alpha\)        & 0.1   \\
    Discount factor       & \(\gamma\)        & 0.99  \\
    Exploration rate      & \(\epsilon_0\)    & 1.0   \\
    Episodes              & \(N\)             & 2000  \\
    \bottomrule
  \end{tabular}
\end{table}

\subsection{Results}

Figure~\ref{fig:learning_curve} shows the learning curve averaged over
10 independent runs.
The agent achieves a mean episode reward of \(XXX \pm YYY\) after \(N\)
training episodes, compared to \(XXX \pm YYY\) for a random baseline.

\begin{figure}[t]
  \centering
  \includegraphics[width=\columnwidth]{figures/learning_curve.pdf}
  \caption{Learning curve (mean $\pm$ std over 10 runs).
           The shaded region denotes one standard deviation.}
  \label{fig:learning_curve}
\end{figure}

\subsection{Discussion}

\textit{[Discuss observations, ablations, and comparisons here.]}

% -------------------------------------------------------
% V. Conclusion
% -------------------------------------------------------
\section{Conclusion}
\label{sec:conclusion}

We demonstrated that a topology-preserving 1-D SOM can effectively discretize
the cart-pole state space and enable reliable policy learning with a simple
tabular Q-learner.
The resulting agent is lightweight, interpretable, and competitive with
neural-network baselines on this benchmark.

Future work will investigate 2-D SOM topologies, continuous action spaces,
and transfer to more complex control tasks.

% -------------------------------------------------------
% Acknowledgment
% -------------------------------------------------------
\section*{Acknowledgment}

\textit{[Optional: acknowledge funding sources, computing resources, etc.]}

% -------------------------------------------------------
% References
% -------------------------------------------------------
\bibliographystyle{IEEEtran}
\bibliography{references}

\end{document}
